\textbf{Phần 1. Trắc nghiệm lựa chọn}
\setcounterpageref{ex}{0}

\begin{ex} %[5H5N25]
Hình tam giác có ba góc đều là góc nhọn được gọi là
\choice
{\True Hình tam giác nhọn}
{Hình tam giác vuông}
{Hình tam giác tù}
{Hình tam giác đều}
\loigiai{
Hình tam giác có $3$ góc nhọn gọi là hình tam giác nhọn.\\
\textbf{Đáp án của câu hỏi là A}
}
\end{ex}

\begin{ex} %[5H5N25]
Công thức tính diện tích $S$ của hình tam giác có độ dài đáy là $a$ và chiều cao tương ứng là $h$ (cùng đơn vị đo) là
\choice
{$S = a \times h$}
{$S = (a + h) \times 2$}
{$S = \dfrac{a + h}{2}$}
{\True $S = \dfrac{a \times h}{2}$}
\loigiai{
Diện tích hình tam giác bằng độ dài đáy nhân với chiều cao (cùng đơn vị đo) rồi chia cho $2$.\\
\textbf{Đáp án của câu hỏi là D}
}
\end{ex}

\begin{ex} %[5H5H25]
Diện tích của hình tam giác có độ dài đáy $10\text{ cm}$ và chiều cao tương ứng $8\text{ cm}$ là
\choice
{$80\text{ cm}^2$}
{$80\text{ cm}$}
{\True $40\text{ cm}^2$}
{$40\text{ cm}$}
\loigiai{
Diện tích tam giác là: $\dfrac{10 \times 8}{2} = 40\text{ (cm}^2\text{)}$.\\
\textbf{Đáp án của câu hỏi là C}
}
\end{ex}

\begin{ex} %[5H5H25]
Diện tích của hình tam giác có độ dài đáy $5\text{ dm}$ và chiều cao $8\text{ dm}$ là
\choice
{$40\text{ dm}^2$}
{\True $20\text{ dm}^2$}
{$13\text{ dm}^2$}
{$20\text{ m}^2$}
\loigiai{
Diện tích tam giác là: $\dfrac{5 \times 8}{2} = 20\text{ (dm}^2\text{)}$.\\
\textbf{Đáp án của câu hỏi là B}
}
\end{ex}

\begin{ex} %[5H5H25]
Một tấm kính có dạng hình tam giác vuông với hai cạnh góc vuông đều bằng $6\text{ m}$. Diện tích của tấm kính đó là
\choice
{$36\text{ m}^2$}
{$12\text{ m}^2$}
{\True $18\text{ m}^2$}
{$24\text{ m}^2$}
\loigiai{
Diện tích tấm kính là: $\dfrac{6 \times 6}{2} = 18\text{ (m}^2\text{)}$.\\
\textbf{Đáp án của câu hỏi là C}
}
\end{ex}

\begin{ex} %[5H5H25]
Diện tích của hình tam giác có độ dài đáy $2\text{ dm}$ và chiều cao $20\text{ cm}$ là
\choice
{$20\text{ dm}^2$}
{\True $2\text{ dm}^2$}
{$20\text{ cm}^2$}
{$200\text{ dm}^2$}
\loigiai{
Đổi $2\text{ dm} = 20\text{ cm}$. Diện tích là: $\dfrac{20 \times 20}{2} = 200\text{ (cm}^2\text{)} = 2\text{ dm}^2$.\\
\textbf{Đáp án của câu hỏi là B}
}
\end{ex}

\begin{ex} %[5H5H25]
Một hình tam giác có diện tích là $15\text{ cm}^2$ và chiều cao $5\text{ cm}$. Độ dài đáy tương ứng của hình tam giác đó là
\choice
{\True $6\text{ cm}$}
{$3\text{ cm}$}
{$15\text{ cm}$}
{$30\text{ cm}$}
\loigiai{
Độ dài đáy là: $15 \times 2 : 5 = 6\text{ (cm)}$.\\
\textbf{Đáp án của câu hỏi là A}
}
\end{ex}

\begin{ex} %[5H5H25]
Một hình tam giác có diện tích $8\text{ m}^2$ và độ dài đáy $4\text{ m}$. Chiều cao tương ứng của hình tam giác đó là
\choice
{$2\text{ m}$}
{$8\text{ m}$}
{$16\text{ m}$}
{\True $4\text{ m}$}
\loigiai{
Chiều cao của hình tam giác là: $8 \times 2 : 4 = 4\text{ (m)}$.\\
\textbf{Đáp án của câu hỏi là D}
}
\end{ex}

\begin{ex} %[5H5V25]
Tờ giấy hình vuông $ABCD$ có cạnh $12\text{ cm}$. Gọi $M$ là trung điểm của $BC$, $N$ là trung điểm của $CD$. Cắt lấy tam giác vuông $MNC$ làm đuôi con cá. Diện tích hình tam giác $MNC$ bằng
\choice
{$36\text{ cm}^2$}
{\True $18\text{ cm}^2$}
{$27\text{ cm}^2$}
{$72\text{ cm}^2$}
\loigiai{
Độ dài đoạn $MC = NC = 12 : 2 = 6\text{ (cm)}$.\\
Diện tích tam giác $MNC$ là: $\dfrac{6 \times 6}{2} = 18\text{ (cm}^2\text{)}$.\\
\textbf{Đáp án của câu hỏi là B}
}
\end{ex}

\begin{ex} %[5H5V25]
Một thửa ruộng hình tam giác có độ dài đáy $32\text{ m}$ và chiều cao $25\text{ m}$. Cứ mỗi $100\text{ m}^2$ người ta thu hoạch được $70\text{ kg}$ thóc. Số thóc thu hoạch được trên thửa ruộng đó là
\choice
{\True $280\text{ kg}$}
{$28\text{ kg}$}
{$560\text{ kg}$}
{$400\text{ kg}$}
\loigiai{
Diện tích thửa ruộng là: $\dfrac{32 \times 25}{2} = 400\text{ (m}^2\text{)}$.\\
Số thóc thu hoạch được là: $400 : 100 \times 70 = 280\text{ (kg)}$.\\
\textbf{Đáp án của câu hỏi là A}
}
\end{ex}

\begin{ex} %[5H5V25]
Nếu tăng độ dài đáy của một hình tam giác lên $2$ lần và giữ nguyên chiều cao thì diện tích của hình tam giác đó sẽ
\choice
{Không thay đổi}
{Tăng lên $4$ lần}
{Giảm đi $2$ lần}
{\True Tăng lên $2$ lần}
\loigiai{
Diện tích mới: $S' = \dfrac{(2 \times a) \times h}{2} = 2 \times \dfrac{a \times h}{2} = 2 \times S$. Vậy diện tích tăng lên $2$ lần.\\
\textbf{Đáp án của câu hỏi là D}
}
\end{ex}

\begin{ex} %[5H5V25]
Một lá cờ hình tam giác vuông có hai cạnh góc vuông lần lượt là $20\text{ cm}$ và $30\text{ cm}$. Người ta cắt $50$ lá cờ như vậy từ một tấm vải. Tổng diện tích vải của $50$ lá cờ đó theo đơn vị mét vuông là
\choice
{$15\text{ m}^2$}
{$0{,}15\text{ m}^2$}
{\True $1{,}5\text{ m}^2$}
{$150\text{ m}^2$}
\loigiai{
Diện tích $1$ lá cờ là: $\dfrac{20 \times 30}{2} = 300\text{ (cm}^2\text{)}$.\\
Diện tích $50$ lá cờ là: $50 \times 300 = 15\,000\text{ (cm}^2\text{)} = 1{,}5\text{ m}^2$.\\
\textbf{Đáp án của câu hỏi là C}
}
\end{ex}

\textbf{Phần 2. Trắc nghiệm đúng sai}
\setcounterpageref{ex}{0}

\begin{ex} %[5H5H25]
Xét tính đúng sai của các khẳng định sau về hình tam giác:
\choiceTF
{\True Hình tam giác đều có $3$ cạnh bằng nhau và $3$ góc đều bằng $60^\circ$.}
{Trong hình tam giác vuông, hai cạnh góc vuông không thể coi là chiều cao và cạnh đáy.}
{\True Diện tích tam giác có đáy $4\text{ cm}$, chiều cao $3\text{ cm}$ là $6\text{ cm}^2$.}
{Hình tam giác có $1$ góc tù và $2$ góc nhọn gọi là hình tam giác vuông.}
\loigiai{
\begin{itemize}
    \item Hình tam giác đều có 3 cạnh bằng nhau và 3 góc bằng nhau cùng bằng $60^\circ$.
    \item Trong tam giác vuông, hai cạnh góc vuông chính là đáy và chiều cao tương ứng của nhau (khẳng định sai).
    \item $S = \dfrac{4 \times 3}{2} = 6\text{ cm}^2$.
    \item Hình tam giác có 1 góc tù gọi là tam giác tù, không phải tam giác vuông.
\end{itemize}
\begin{itemize}
    \item a) \textbf{Đúng}.
    \item b) \textbf{Sai}.
    \item c) \textbf{Đúng}.
    \item d) \textbf{Sai}.
\end{itemize}
}
\end{ex}

\begin{ex} %[5H5V25]
Một hình tam giác $ABC$ có đáy $BC = 10\text{ dm}$, chiều cao $AH = 6\text{ dm}$. Xét tính đúng sai của các khẳng định sau:
\choiceTF
{\True Diện tích hình tam giác $ABC$ là $30\text{ dm}^2$.}
{Đổi diện tích hình tam giác ra mét vuông là $3\text{ m}^2$.}
{\True Nếu kéo dài đáy $BC$ thêm $2\text{ dm}$ thì diện tích tăng thêm $6\text{ dm}^2$.}
{\True Tỉ số diện tích phần tăng thêm và diện tích ban đầu là $\dfrac{1}{5}$.}
\loigiai{
\begin{itemize}
    \item Diện tích tam giác: $\dfrac{10 \times 6}{2} = 30\text{ dm}^2$.
    \item $30\text{ dm}^2 = 0{,}3\text{ m}^2 \neq 3\text{ m}^2$.
    \item Diện tích tăng thêm: $\dfrac{2 \times 6}{2} = 6\text{ dm}^2$.
    \item Tỉ số: $6 : 30 = \dfrac{1}{5}$.
\end{itemize}
\begin{itemize}
    \item a) \textbf{Đúng}.
    \item b) \textbf{Sai}: Kết quả là $0{,}3\text{ m}^2$.
    \item c) \textbf{Đúng}.
    \item d) \textbf{Đúng}.
\end{itemize}
}
\end{ex}

\textbf{Phần 3. Trả lời ngắn}
\setcounterpageref{ex}{0}

\begin{ex} %[5H5H25]
Tính diện tích hình tam giác có đáy $4\text{ cm}$ và chiều cao $3\text{ cm}$.
\shortans{6}
\loigiai{
\textbf{Đáp số:} $6$\\
$\dfrac{4 \times 3}{2} = 6\text{ (cm}^2\text{)}$.
}
\end{ex}

\begin{ex} %[5H5H25]
Tính diện tích tấm kính hình tam giác vuông có hai cạnh góc vuông cùng bằng $6\text{ m}$.
\shortans{18}
\loigiai{
\textbf{Đáp số:} $18$\\
$\dfrac{6 \times 6}{2} = 18\text{ (m}^2\text{)}$.
}
\end{ex}

\begin{ex} %[5H5H25]
Hình tam giác có diện tích $25\text{ dm}^2$ và đáy $10\text{ dm}$. Tính chiều cao tương ứng của hình tam giác đó theo đơn vị đề-xi-mét.
\shortans{5}
\loigiai{
\textbf{Đáp số:} $5$\\
Chiều cao là: $25 \times 2 : 10 = 5\text{ (dm)}$.
}
\end{ex}

\begin{ex} %[5H5V25]
Tờ giấy hình vuông cạnh $12\text{ cm}$ cắt bớt tam giác vuông ở góc có hai cạnh góc vuông dài $6\text{ cm}$ và $6\text{ cm}$. Tính diện tích phần tam giác bị cắt theo đơn vị xăng-ti-mét vuông.
\shortans{18}
\loigiai{
\textbf{Đáp số:} $18$\\
$\dfrac{6 \times 6}{2} = 18\text{ (cm}^2\text{)}$.
}
\end{ex}

\begin{ex} %[5H5V25]
Một thửa ruộng hình tam giác có đáy $32\text{ m}$, chiều cao $25\text{ m}$. Tính diện tích thửa ruộng đó theo đơn vị mét vuông.
\shortans{400}
\loigiai{
\textbf{Đáp số:} $400$\\
$\dfrac{32 \times 25}{2} = 400\text{ (m}^2\text{)}$.
}
\end{ex}

\begin{ex} %[5H5V25]
Cắt $50$ lá cờ hình tam giác vuông có kích thước hai cạnh góc vuông là $20\text{ cm}$ và $30\text{ cm}$. Hỏi tổng diện tích vải của $50$ lá cờ đó là bao nhiêu mét vuông? (Nhập kết quả dưới dạng số thập phân).
\shortans{1,5}
\loigiai{
\textbf{Đáp số:} $1{,}5$\\
$50 \times \dfrac{0{,}2 \times 0{,}3}{2} = 1{,}5\text{ (m}^2\text{)}$.
}
\end{ex}
